\chapter{Normal Reference Ranges --- What the Numbers Mean}
\label{app:reference_ranges}

\clearpage
\section{Introduction}

Every lab result comes with a ``reference range'' --- the set of values considered normal for a healthy person. But ``normal'' is not one-size-fits-all. Your age, sex, medications, time of day, and even whether you fasted all affect your numbers.

This appendix gives you the reference ranges for the most common blood tests, organized by body system, with plain-language explanations of what each number tells your doctor and what it means for you.

A few critical points before you start reading:
\begin{itemize}
    \item Reference ranges vary by lab. The numbers in this appendix are typical, but your lab may use slightly different cutoffs based on the equipment they use and the population they serve.
    \item Being slightly outside the reference range is not necessarily a problem. One abnormal result rarely means you are sick. Your doctor interprets these numbers in the context of your symptoms, medical history, and other test results.
    \item Some tests have different ranges for men and women. Where applicable, both ranges are listed.
\end{itemize}

\clearpage
\section{Hematology --- Your Blood Count}

A complete blood count (CBC) is one of the most commonly ordered tests. It tells your doctor about your red blood cells (which carry oxygen), white blood cells (which fight infection), and platelets (which help your blood clot).

\subsection{Understanding Your Results}

\begin{longtable}{@{}llll@{}}
\toprule
\textbf{Test} & \textbf{Male} & \textbf{Female} & \textbf{Units} \\
\midrule
\endhead
Hemoglobin (Hgb) & 13.5--17.5 & 12.0--15.5 & g/dL \\
Hematocrit (Hct) & 41--53 & 36--46 & \% \\
Red Blood Cell Count (RBC) & 4.5--5.9 & 4.0--5.2 & $\times 10^{12}$/L \\
Mean Corpuscular Volume (MCV) & 80--100 & 80--100 & fL \\
Mean Corpuscular Hemoglobin (MCH) & 27--33 & 27--33 & pg \\
Mean Corpuscular Hemoglobin Concentration (MCHC) & 32--36 & 32--36 & g/dL \\
Red Cell Distribution Width (RDW) & 11.5--14.5 & 11.5--14.5 & \% \\
Reticulocyte Count & 0.5--1.5 & 0.5--1.5 & \% \\
\bottomrule
\end{longtable}

\subsubsection{What These Numbers Mean For You}

\textbf{Hemoglobin and Hematocrit.} Hemoglobin is the protein inside your red blood cells that grabs oxygen from your lungs and delivers it to every tissue in your body. Hematocrit is the percentage of your blood volume that is made up of red blood cells. When these are low, you are anemic --- you may feel tired, short of breath, dizzy, or look pale. When these are high, it could mean you are dehydrated, you smoke, or in rare cases, your bone marrow is making too many red blood cells.

\textbf{MCV --- the size of your red blood cells.} This number is critical because it helps your doctor figure out \emph{why} your hemoglobin is abnormal. Small red blood cells (low MCV) usually point to iron deficiency or thalassemia. Large red blood cells (high MCV) usually point to vitamin B12 or folate deficiency, liver disease, or hypothyroidism. Normal-sized red blood cells with a low hemoglobin could mean anemia of chronic disease, kidney disease, or acute blood loss.

\textbf{RDW.} This measures how much variation there is in the size of your red blood cells. A high RDW means your red blood cells are different sizes --- this is often an early sign of iron deficiency, even before your hemoglobin drops.

\textbf{Reticulocyte Count.} These are immature red blood cells freshly released from your bone marrow. A high reticulocyte count means your bone marrow is responding to a problem (like bleeding or hemolysis) by working overtime. A low count means your bone marrow is not responding --- it may be suppressed by disease, medication, or nutrient deficiency.

\subsection{Understanding Your Results}

\begin{longtable}{@{}lll@{}}
\toprule
\textbf{Test} & \textbf{Reference Range} & \textbf{Units} \\
\midrule
Total White Blood Cell Count (WBC) & 4.5--11.0 & $\times 10^9$/L \\
Neutrophils & 40--70 & \% \\
Lymphocytes & 20--40 & \% \\
Monocytes & 2--8 & \% \\
Eosinophils & 1--4 & \% \\
Basophils & 0--1 & \% \\
\bottomrule
\end{longtable}

\subsubsection{What These Numbers Mean For You}

\textbf{Total WBC.} White blood cells are your immune army. A high WBC usually means your body is fighting an infection, but it can also be caused by stress, steroids, inflammation, or in rare cases, leukemia. A low WBC (called leukopenia) means your immune system may be weakened --- this can happen from viral infections, autoimmune diseases, bone marrow problems, or as a side effect of medications like chemotherapy.

\textbf{The differential (the percentage of each type of white blood cell).} Neutrophils are your first responders against bacterial infections. If your neutrophils are elevated, your doctor will suspect a bacterial infection. Lymphocytes fight viral infections and are important for your adaptive immune system. Elevated eosinophils are a red flag for allergies or parasitic infections. Each type of white blood cell rising or falling tells a different story about what is happening inside your body.

\subsection{Understanding Your Results}

\begin{longtable}{@{}lll@{}}
\toprule
\textbf{Test} & \textbf{Reference Range} & \textbf{Units} \\
\midrule
Platelet Count & 150--450 & $\times 10^9$/L \\
Mean Platelet Volume (MPV) & 7.4--10.4 & fL \\
\bottomrule
\end{longtable}

\subsubsection{What These Numbers Mean For You}

\textbf{Platelets.} These are the tiny cell fragments that plug cuts and stop bleeding. Low platelets (thrombocytopenia) mean you are at risk for bruising and bleeding --- causes include immune destruction, liver disease, infections, and certain medications. High platelets (thrombocytosis) can mean inflammation, iron deficiency, or a bone marrow disorder.

\textbf{MPV} tells your doctor the average size of your platelets. Large platelets are often freshly produced, meaning your bone marrow is actively making them. Small platelets may be older platelets that have been circulating for a while.

\clearpage
\section{Basic Metabolic Panel --- Your Chemistry Check}

The basic metabolic panel (BMP) is a snapshot of your body's chemical balance. It checks your blood sugar, kidney function, and electrolytes --- the minerals that control everything from your heartbeat to your muscle contractions.

\begin{longtable}{@{}lll@{}}
\toprule
\textbf{Analyte} & \textbf{Reference Range} & \textbf{Units} \\
\midrule
Sodium (Na) & 136--145 & mmol/L \\
Potassium (K) & 3.5--5.1 & mmol/L \\
Chloride (Cl) & 98--107 & mmol/L \\
Bicarbonate (CO$_2$) & 22--29 & mmol/L \\
Blood Urea Nitrogen (BUN) & 7--20 & mg/dL \\
Creatinine & 0.7--1.3 (M), 0.6--1.1 (F) & mg/dL \\
Glucose (fasting) & 70--99 & mg/dL \\
Calcium (Ca) & 8.5--10.5 & mg/dL \\
\bottomrule
\end{longtable}

\subsubsection{What These Numbers Mean For You}

\textbf{Sodium.} This is your body's most important electrolyte for maintaining fluid balance. Low sodium (hyponatremia) can cause confusion, seizures, and in severe cases, coma. It is commonly caused by drinking too much water, heart failure, kidney disease, or certain medications. High sodium (hypernatremia) usually means dehydration.

\textbf{Potassium.} This electrolyte controls your heartbeat. Even small abnormalities can cause dangerous arrhythmias. Low potassium (hypokalemia) commonly results from vomiting, diarrhea, or diuretic medications. High potassium (hyperkalemia) is a medical emergency that can cause cardiac arrest --- it is seen in kidney failure, certain medications (like ACE inhibitors and potassium-sparing diuretics), and tissue destruction (burns, crush injuries).

\textbf{BUN and Creatinine --- kidney function.} These are waste products filtered out by your kidneys. When they are elevated, it means your kidneys are not filtering well enough. BUN can also rise with dehydration and high-protein diets. Creatinine is a more reliable indicator of kidney function than BUN. Your doctor also calculates your eGFR (estimated glomerular filtration rate) from your creatinine, which tells you what percentage of your kidney function remains.

\textbf{Glucose (fasting).} This is your blood sugar after not eating for at least 8 hours. A fasting glucose of 100--125 mg/dL means you have prediabetes. A fasting glucose of 126 mg/dL or higher, confirmed on two separate occasions, means you have diabetes. A fasting glucose below 70 mg/dL is hypoglycemia, which can cause shakiness, sweating, confusion, and in severe cases, loss of consciousness.

\textbf{Calcium.} Calcium is not just for bones --- it is essential for muscle contraction, nerve signaling, and blood clotting. High calcium can be a sign of overactive parathyroid glands or certain cancers. Low calcium can cause tingling, muscle cramps, and seizures.

\clearpage
\section{Liver Function Tests --- Is Your Liver Working Properly?}

Your liver performs over 500 functions, including filtering toxins, producing proteins, and making bile for digestion. Liver function tests (LFTs) measure how well your liver is doing its job.

\begin{longtable}{@{}lll@{}}
\toprule
\textbf{Analyte} & \textbf{Reference Range} & \textbf{Units} \\
\midrule
Total Bilirubin & 0.3--1.2 & mg/dL \\
Direct Bilirubin & 0.0--0.3 & mg/dL \\
Alkaline Phosphatase (ALP) & 40--129 & U/L \\
ALT & 7--56 (M), 7--45 (F) & U/L \\
AST & 13--39 (M), 13--35 (F) & U/L \\
GGT & 9--48 (M), 9--32 (F) & U/L \\
Total Protein & 6.0--8.3 & g/dL \\
Albumin & 3.5--5.0 & g/dL \\
\bottomrule
\end{longtable}

\subsubsection{What These Numbers Mean For You}

\textbf{ALT and AST --- liver cell damage.} ALT (alanine aminotransferase) and AST (aspartate aminotransferase) are enzymes that leak out of damaged liver cells. ALT is more specific to the liver, while AST is also found in the heart and muscles. Elevations can mean hepatitis (viral, alcoholic, or drug-induced), fatty liver disease, or liver injury from medications like acetaminophen. Very high elevations (over 1000 U/L) are seen in acute liver failure from toxins, viral hepatitis, or ischemia (lack of blood flow to the liver).

\textbf{ALP and GGT --- bile duct problems.} ALP (alkaline phosphatase) is elevated when bile flow is blocked or when bone is being remodeled. GGT (gamma-glutamyl transferase) helps distinguish whether an elevated ALP is coming from the liver or the bones --- if both ALP and GGT are high, the problem is in the liver. Bile duct obstruction from gallstones, tumors, or scarring will raise both of these.

\textbf{Bilirubin --- jaundice.} Bilirubin is a yellow pigment produced when old red blood cells break down. The liver processes it and puts it into bile. When bilirubin is high, your skin and eyes turn yellow (jaundice). High bilirubin can mean liver disease, bile duct obstruction, or excessive red blood cell destruction (hemolysis).

\textbf{Albumin and Total Protein.} Albumin is a protein made exclusively by your liver. Low albumin means your liver is not producing enough protein, or you are losing it through your kidneys (nephrotic syndrome) or your gut (protein-losing enteropathy). Low albumin causes swelling (edema) in your legs and abdomen.

\clearpage
\section{Thyroid Function --- Is Your Metabolism Regulated?}

Your thyroid gland produces hormones that control your metabolism --- how fast your body uses energy. Thyroid dysfunction is extremely common and often underdiagnosed.

\begin{longtable}{@{}lll@{}}
\toprule
\textbf{Analyte} & \textbf{Reference Range} & \textbf{Units} \\
\midrule
TSH & 0.4--4.0 & mIU/L \\
Free T4 & 0.8--1.8 & ng/dL \\
Free T3 & 2.0--4.4 & pg/mL \\
\bottomrule
\end{longtable}

\subsubsection{What These Numbers Mean For You}

\textbf{TSH --- the master switch.} TSH (thyroid-stimulating hormone) is produced by your pituitary gland in the brain. When your thyroid is underactive, TSH goes up (your brain is screaming at the thyroid to work harder). When your thyroid is overactive, TSH goes down (your brain is trying to shut it down). TSH is the single best screening test for thyroid disease.

\textbf{Free T4 and Free T3 --- the actual thyroid hormones.} Free T4 is the main hormone your thyroid produces. Free T3 is the more active form, converted from T4 in your tissues. These are measured as ``free'' hormone (unbound to proteins) because that is the biologically active portion.

\textbf{Interpreting the pattern:} High TSH with low Free T4 means overt hypothyroidism --- your thyroid is failing. High TSH with normal Free T4 means subclinical hypothyroidism --- your thyroid is struggling but still compensating. Low TSH with high Free T4 means overt hyperthyroidism. Low TSH with normal Free T4 means subclinical hyperthyroidism or a pituitary problem.

\clearpage
\section{Coagulation --- How Well Does Your Blood Clot?}

Coagulation tests measure how well your blood can form clots. These tests are essential before surgery, for monitoring patients on blood thinners, and for diagnosing bleeding disorders.

\begin{longtable}{@{}lll@{}}
\toprule
\textbf{Test} & \textbf{Reference Range} & \textbf{Units} \\
\midrule
Prothrombin Time (PT) & 11--13.5 & seconds \\
INR & 0.8--1.2 & ratio \\
aPTT & 25--35 & seconds \\
Fibrinogen & 200--400 & mg/dL \\
D-dimer & $<$0.5 & $\mu$g/mL FEU \\
\bottomrule
\end{longtable}

\subsubsection{What These Numbers Mean For You}

\textbf{PT and INR.} The prothrombin time measures how long it takes your blood to clot through the extrinsic pathway. The INR (international normalized ratio) is a standardized version of the PT --- it allows results to be compared across different labs. If you are on warfarin (Coumadin), your doctor wants your INR to be in a therapeutic range (usually 2.0--3.0) --- high enough to prevent clots, but not so high that you bleed uncontrollably. An INR above 1.2 in someone not on blood thinners may indicate liver disease or a clotting factor deficiency.

\textbf{aPTT.} This measures the intrinsic clotting pathway. It is used to monitor heparin therapy. A prolonged aPTT in someone not on heparin may suggest liver disease, hemophilia, or the presence of lupus anticoagulant.

\textbf{Fibrinogen.} This is the protein that is converted into fibrin to form the actual clot structure. Low fibrinogen means your blood cannot form stable clots --- this is seen in severe liver disease, disseminated intravascular coagulation (DIC), and massive bleeding.

\textbf{D-dimer.} When a clot forms in your body and then breaks down, it releases D-dimer. A negative D-dimer is very good at ruling out blood clots (like deep vein thrombosis and pulmonary embolism) in low-risk patients. A positive D-dimer means clots are being formed and broken down somewhere, but it is not specific --- inflammation, surgery, pregnancy, and cancer can all raise it.

\clearpage
\section{Lipid Panel --- Your Heart Disease Risk}

\begin{longtable}{@{}lll@{}}
\toprule
\textbf{Analyte} & \textbf{Desirable} & \textbf{Units} \\
\midrule
Total Cholesterol & $<$200 & mg/dL \\
LDL Cholesterol & $<$100 & mg/dL \\
HDL Cholesterol & $>$40 (M), $>$50 (F) & mg/dL \\
Triglycerides & $<$150 & mg/dL \\
\bottomrule
\end{longtable}

\subsubsection{What These Numbers Mean For You}

\textbf{Total Cholesterol.} This is the sum of all cholesterol in your blood. Below 200 is desirable. 200--239 is borderline high. 240 and above is high.

\textbf{LDL --- the ``bad'' cholesterol.} LDL (low-density lipoprotein) carries cholesterol into your artery walls, where it forms plaques that cause heart attacks and strokes. Lower is better. Below 100 is optimal. If you already have heart disease, your doctor may want your LDL below 70.

\textbf{HDL --- the ``good'' cholesterol.} HDL (high-density lipoprotein) carries cholesterol away from your arteries back to the liver for disposal. Higher is better. Below 40 for men and below 50 for women puts you at increased risk for heart disease.

\textbf{Triglycerides.} These are a type of fat in your blood. They rise with a high-sugar diet, excessive alcohol, obesity, and uncontrolled diabetes. Very high triglycerides (over 500 mg/dL) increase your risk of pancreatitis.

\clearpage
\section{Renal Function --- Are Your Kidneys Healthy?}

\begin{longtable}{@{}lll@{}}
\toprule
\textbf{Analyte} & \textbf{Reference Range} & \textbf{Units} \\
\midrule
BUN & 7--20 & mg/dL \\
Creatinine & 0.7--1.3 (M), 0.6--1.1 (F) & mg/dL \\
eGFR & $>$90 & mL/min/1.73 m$^2$ \\
Uric Acid & 3.4--7.0 (M), 2.4--6.0 (F) & mg/dL \\
\bottomrule
\end{longtable}

\subsubsection{What These Numbers Mean For You}

\textbf{eGFR --- the most important kidney number.} This is calculated from your creatinine, age, and sex. It estimates how well your kidneys are filtering waste. A value above 90 is normal. 60--89 means mild kidney damage. 30--59 means moderate kidney failure. Below 30 means severe kidney failure. Below 15 means kidney failure requiring dialysis or transplant.

\textbf{Uric Acid.} High uric acid is the cause of gout --- the excruciatingly painful joint inflammation, usually in the big toe. It can also form kidney stones. High uric acid is caused by a diet rich in red meat and shellfish, excessive alcohol, and certain medications.

\clearpage
\section{Endocrine --- Diabetes and Hormone Testing}

\begin{longtable}{@{}lll@{}}
\toprule
\textbf{Analyte} & \textbf{Reference Range} & \textbf{Units} \\
\midrule
HbA1c & $<$5.7 & \% \\
Fasting Insulin & 2--25 & $\mu$IU/mL \\
Cortisol (morning) & 5--25 & $\mu$g/dL \\
ACTH & 10--60 & pg/mL \\
\bottomrule
\end{longtable}

\subsubsection{What These Numbers Mean For You}

\textbf{HbA1c --- your 3-month blood sugar average.} Unlike a fasting glucose that gives you one snapshot, HbA1c reflects your average blood sugar over the past 2--3 months by measuring how much sugar is stuck to your hemoglobin. Below 5.7\% is normal. 5.7--6.4\% is prediabetes. 6.5\% or above is diabetes. If you are diabetic, most doctors aim for below 7\%.

\textbf{Cortisol (morning).} Cortisol is your ``stress hormone,'' and it peaks in the morning. Low morning cortisol may suggest adrenal insufficiency (Addison's disease). High cortisol may suggest Cushing's syndrome or chronic stress.

\textbf{ACTH.} This hormone from your pituitary tells your adrenal glands to make cortisol. Low ACTH with low cortisol suggests a pituitary problem. High ACTH with low cortisol suggests an adrenal problem (Addison's disease).

\clearpage
\section{Arterial Blood Gas --- How Well Are You Breathing?}

An arterial blood gas (ABG) is a needle stick directly into an artery (usually in the wrist) that measures how well your lungs are oxygenating your blood and how well your body is managing acid-base balance.

\begin{longtable}{@{}lll@{}}
\toprule
\textbf{Parameter} & \textbf{Reference Range} & \textbf{Units} \\
\midrule
pH & 7.35--7.45 & -- \\
PaCO$_2$ & 35--45 & mmHg \\
PaO$_2$ & 80--100 & mmHg \\
HCO$_3$ & 22--26 & mmol/L \\
Base Excess & $-2$ to $+2$ & mmol/L \\
SaO$_2$ & 95--100 & \% \\
\bottomrule
\end{longtable}

\subsubsection{What These Numbers Mean For You}

\textbf{pH.} Your blood must stay between 7.35 and 7.45 to keep your enzymes, cells, and organs functioning. Below 7.35 is acidosis --- too much acid in your blood. This can be caused by kidney failure, diabetic ketoacidosis, or lung failure retaining CO$_2$. Above 7.45 is alkalosis --- too little acid. This can be caused by vomiting, hyperventilation, or excessive antacid use.

\textbf{PaCO$_2$ (partial pressure of carbon dioxide).} This reflects how well your lungs are blowing off CO$_2$. High PaCO$_2$ means your lungs are failing (respiratory acidosis). Low PaCO$_2$ means you are hyperventilating (respiratory alkalosis).

\textbf{PaO$_2$ (partial pressure of oxygen).} This is the amount of oxygen dissolved in your arterial blood. Below 80 mmHg is hypoxemia --- not enough oxygen is getting into your blood. This can be from pneumonia, pulmonary embolism, heart failure, or lung disease.

\textbf{HCO$_3$ (bicarbonate).} This is your body's buffer system. When your blood is too acidic, your kidneys compensate by holding onto bicarbonate. When your blood is too alkaline, your kidneys excrete bicarbonate. The relationship between PaCO$_2$ and HCO$_3$ tells your doctor whether the problem is primarily respiratory (lung-driven) or metabolic (kidney-driven).

\textbf{SaO$_2$ (oxygen saturation).} This is the percentage of your hemoglobin that is carrying oxygen. Below 90\% is dangerously low. A normal pulse oximeter on your finger measures this same number non-invasively.

\clearpage
\section{How to Use This Information}

Reference ranges are a starting point, not a diagnosis. Here is what you should actually do with this information:

\begin{enumerate}
    \item \textbf{Do not panic over one abnormal result.} Many things can temporarily shift your numbers --- dehydration, a recent meal, exercise, stress, or medications. Your doctor will repeat abnormal results before making any diagnosis.

    \item \textbf{Compare your results over time.} A trend is more meaningful than a single value. A creatinine that has been creeping up over six months is more concerning than one slightly elevated result.

    \item \textbf{Ask your doctor what the numbers mean in your context.} A hemoglobin of 13.2 g/dL is ``normal'' for a man but would be high for a pregnant woman. An INR of 1.5 is fine if you are not on blood thinners but dangerous if you are supposed to be on warfarin.

    \item \textbf{Know your own baseline.} If you have periodic blood work done, keep a record of your results. Knowing your personal normal range is more useful than any textbook range.

    \item \textbf{Bring this appendix to your next appointment.} Reference ranges are not secrets --- you have every right to understand your own test results.
\end{enumerate}
